% ===========================================================================
\section{Uncertainty and physical controls}
\label{sec:uncertainty-controls}
% ===========================================================================

\subsection{Parametric uncertainty}

The legacy experiment samples each published 2006 factor from a bounded log-uniform
distribution and reports the median, P2.5, P97.5 and full methodological bounds. Rows in the
same physical stream share their range quantile.

For the 2019 inventory, published IPCC limits are interpreted as bounded expert ranges. A
split-triangular distribution places the reported value at the mode and median while preserving
the stated lower and upper limits. One production-volume deviate is shared across gas
production, gathering and processing; oil production and its transport fallback share the oil
deviate; and gases within a physical factor group share a quantile. Additional processing and
coverage residuals remain independent. The simulation uses \num{10000} iterations and seed 42.
The exact transformation, correlation groups and material inputs are reported in the supplement.

Input importance is measured by Spearman rank correlation with the national total. Normalized
squared correlations provide a screening measure of the parameters controlling the conditional
interval. A distribution-shape sensitivity reruns the reference model with bounded uniform and
bounded log-uniform inputs.

\subsection{Structural configurations}

The factorial experiment crosses upstream practice (higher/lower), network practice
(higher/lower), T\&D activity basis (pipeline length/throughput), post-meter industrial coverage
(fertilizer feedstock excluded, the reference/feedstock included/whole fertilizer class
excluded/half/excluded), processing activity (gross
production/marketable supply/excluded) and
sour-gas share (32\%/16\%/excluded). Removing invalid combinations leaves 280 deterministic,
unweighted configurations. No probability weights are assigned. We report one envelope for
models retaining every public-data
segment and a wider envelope that permits proxy-supported processing and post-meter coverage to
be excluded. A separate named \emph{central} configuration takes every stream with published
high/low practice bounds at their geometric mean; it is priced, never selected, and is not a
factorial leg because the practice bounds are readings of one control, not coverage choices.

Parametric intervals are calculated separately for five interpretable configurations:
reference, central, lower practices, no post-meter industrial/power, and the
minimum-coverage combined
case. This keeps uncertainty in measured activities and factors conditional on the selected
infrastructure model.

\subsection{Gas-flow balance and UFG diagnostic}

For each matched year, the physical account separates upstream and downstream identities:
\begin{align}
F_{gross}&=F_{indigenous,marketable}+R_{upstream},\\
F_{indigenous,marketable}+F_{RLNG}&=F_{use}+F_{UFG}+R_{downstream}.
\end{align}
The residuals preserve differences among published boundaries. The upstream disposition table
also records reinjection, own use, processing shrinkage, venting and flaring where public
quantities are available. Calendar-2024 VIIRS flare volume is shown as a volume-equivalent memo
against the adjacent gross-to-marketable difference.

UFG includes physical losses, theft, metering error and accounting effects. It is therefore used
only as a plausibility bound. If fraction $q$ were physical loss, the available methane would be
\begin{equation}
M_{\mathrm{CH_4}}=q V_{\mathrm{UFG}}\rho_{gas}x_{\mathrm{CH_4}}.
\end{equation}
Reference and lower-practice transmission-plus-distribution methane are expressed as the implied
fraction $q$; post-meter emissions are outside this utility balance.
